\section{Code and Data Availability}
\label{sec:code}
All code is written in Python $\geq$3.12, with \texttt{scikit-learn}~\cite{sklearn} for the classical machine-learning models and \texttt{TensorFlow}/\texttt{Keras}~\cite{tensorflow} for the neural networks. Key package versions: \texttt{uproot 5.7}, \texttt{iminuit 2.32.0}, \texttt{emcee 3.1.6}, \texttt{corner 2.3.0}, \texttt{scikit-learn 1.9}, \texttt{tensorflow 2.21.0}, \texttt{keras 3.15.1}, \texttt{scipy 1.18}, \texttt{numpy 2.4}, \texttt{pandas 3.0}.

The dataset is the CERN Open Data record \href{https://opendata.cern.ch/record/4900}{4900} ("Matter Antimatter Differences -- B meson decays to three hadrons"). The analysis code is available in the accompanying repository, in \texttt{Code\_general/Frequentist/} (frequentist fits, the common systematic study, and \texttt{freq\_validation.py} with the pseudo-experiments, profile likelihoods, confidence belt and bootstrap studies of Sec.~\ref{sec:validation} and Appendix~\ref{sec:appendix_checks}), \texttt{Code\_general/Bayesian/} (nominal model, its checks, the minimal model, and \texttt{bayes\_validation.py} and \texttt{bayes\_extras.py} for the checks of the posterior), and \texttt{Code\_general/Machine\_Learning/}; the figures of this paper are produced from the saved results of those analyses by the scripts in \texttt{Report/figures/} (\texttt{plot\_style.py} for the common style, \texttt{regen\_mass\_spectrum.py}, \texttt{regen\_figures\_nominal.py}, \texttt{regen\_figures\_ml\_nominal.py}, \texttt{regen\_figures\_validation.py}). All random numbers are generated with fixed seeds, so the pseudo-experiments can be reproduced.
